% Protocol Specification -- Multi-Client File Transfer
%
% Compile on Overleaf (overleaf.com, free, no install): create a new
% project, paste this file's content as main.tex, click Recompile.
% Or locally with a TeX distribution: pdflatex protocol_spec.tex (run
% twice if section numbers/links look off on the first pass).

\documentclass[11pt]{article}

\usepackage[margin=1in]{geometry}
\usepackage{longtable}
\usepackage{array}
\usepackage{booktabs}
\usepackage{enumitem}
\usepackage{hyperref}
\usepackage{parskip}
\usepackage{titlesec}

\hypersetup{colorlinks=true, linkcolor=blue, urlcolor=blue}
\titleformat{\section}{\large\bfseries}{\thesection}{1em}{}
\titleformat{\subsection}{\normalsize\bfseries}{\thesubsection}{1em}{}

\newcommand{\cmd}[1]{\texttt{#1}}

\title{Application-Layer Protocol Specification\\
\large Multi-Client File Transfer}
\author{}
\date{}

\begin{document}
\maketitle

\section{Overview}

This document specifies a custom, text-based, line-oriented application
protocol running on top of TCP. It is used by the multi-client file
transfer system to handle authentication, presence, private/broadcast
messaging, and file transfer (client$\leftrightarrow$server and
client$\leftrightarrow$client via the server).

The protocol is connection-oriented: a client opens one TCP connection to
the server, authenticates, and keeps that connection open for the rest of
the session (all messaging and file transfer happen over that same
socket).

\section{Message Format}

Every \textbf{control message} is a single line of UTF-8 text, terminated
by \cmd{\textbackslash{}n}:

\begin{center}
\texttt{COMMAND|FIELD1|FIELD2|...}
\end{center}

\begin{itemize}[leftmargin=1.5em]
  \item Fields are separated by the \cmd{|} character.
  \item The final field of a message may itself contain \cmd{|}
        characters (used for free-text message bodies in \cmd{MSG} /
        \cmd{BROADCAST}); every other field must not contain \cmd{|} or a
        literal newline.
  \item Every control message is read using a length-agnostic line
        reader (read byte-by-byte until \cmd{\textbackslash{}n}), so no
        message needs to declare its own length.
\end{itemize}

\textbf{File payloads} are the one exception: raw binary bytes cannot
safely be represented as a delimited text line. Whenever a control
message announces a file transfer, it includes an exact byte count as a
field. Immediately after that control line, exactly that many raw bytes
follow with \textbf{no} delimiter or framing --- the receiver already
knows how many bytes to read.

\section{Commands}

\begin{longtable}{@{}p{2.1cm}p{1.1cm}p{5.2cm}p{6.2cm}@{}}
\toprule
\textbf{Command} & \textbf{Dir.} & \textbf{Format} & \textbf{Purpose} \\
\midrule
\endhead
\cmd{LOGIN} & C$\to$S & \cmd{LOGIN|username|password} & Authenticate. Unknown usernames are auto-registered with the given password. \\
\cmd{LOGOUT} & C$\to$S & \cmd{LOGOUT} & Gracefully end the session; server closes the connection. \\
\cmd{LIST} & C$\to$S & \cmd{LIST} & Request the list of currently online users. \\
\cmd{FILES} & C$\to$S & \cmd{FILES} & Request the list of files currently stored on the server (available for \cmd{DOWNLOAD}). \\
\cmd{ALLUSERS} & C$\to$S & \cmd{ALLUSERS} & Request every registered account (from \cmd{users.txt}), each tagged online/offline -- unlike \cmd{LIST}, this includes users who aren't currently connected. \\
\cmd{MSG} & C$\to$S & \cmd{MSG|username|text} & Send a private message to one online user. \\
\cmd{MSG} & S$\to$C & \cmd{MSG|sender|text} & Delivery of a private message (pushed to the recipient). \\
\cmd{BROADCAST} & C$\to$S & \cmd{BROADCAST|text} & Send a message to every other online user. \\
\cmd{BROADCAST} & S$\to$C & \cmd{BROADCAST|sender|text} & Delivery of a broadcast (pushed to every recipient). \\
\cmd{UPLOAD} & C$\to$S & \cmd{UPLOAD|filename|filesize} & Push a file to server storage. Followed by \cmd{filesize} raw bytes once the server replies \cmd{READY}. \\
\cmd{DOWNLOAD} & C$\to$S & \cmd{DOWNLOAD|filename} & Request a file from server storage. \\
\cmd{SEND} & C$\to$S & \cmd{SEND|targetuser|filename|filesize} & Send a file directly to another online user, relayed live through the server. Followed by \cmd{filesize} raw bytes once the server replies \cmd{READY}. \\
\bottomrule
\end{longtable}

\section{Server Responses}

\begin{longtable}{@{}p{2.4cm}p{5.2cm}p{6.9cm}@{}}
\toprule
\textbf{Response} & \textbf{Format} & \textbf{Meaning} \\
\midrule
\endhead
\cmd{LOGIN\_OK} & \cmd{LOGIN\_OK|username} & Login accepted. \\
\cmd{LOGOUT\_OK} & \cmd{LOGOUT\_OK} & Logout acknowledged; connection will be closed. \\
\cmd{USERS} & \cmd{USERS|name1,name2,...} & Reply to \cmd{LIST}. Comma-separated, possibly empty. \\
\cmd{FILES} & \cmd{FILES|file1,file2,...} & Reply to \cmd{FILES}. Comma-separated, possibly empty. \\
\cmd{ALLUSERS} & \cmd{ALLUSERS|user1:online,...} & Reply to \cmd{ALLUSERS}. Each entry is \cmd{username:online} or \cmd{username:offline}. \\
\cmd{MSG\_OK} & \cmd{MSG\_OK} & Private message was delivered. \\
\cmd{BROADCAST\_OK} & \cmd{BROADCAST\_OK} & Broadcast was relayed to all other online users. \\
\cmd{READY} & \cmd{READY} & Server is ready to receive the raw file bytes for an \cmd{UPLOAD} or \cmd{SEND}. \\
\cmd{UPLOAD\_OK} & \cmd{UPLOAD\_OK|filename} & File stored successfully. \\
\cmd{DOWNLOAD\_OK} & \cmd{DOWNLOAD\_OK|filename|filesize} & Followed immediately by \cmd{filesize} raw bytes. \\
\cmd{SEND\_OK} & \cmd{SEND\_OK|filename} & File was relayed to the target user. \\
\cmd{SEND\_INCOMING} & \cmd{SEND\_INCOMING|filename|sender|filesize} & \textbf{Push}, not a reply: notifies the recipient of an inbound \cmd{SEND}. Followed immediately by \cmd{filesize} raw bytes. \\
\cmd{ERROR} & \cmd{ERROR|CODE} & See Section~5. \\
\bottomrule
\end{longtable}

Note the asymmetry: \cmd{MSG}, \cmd{BROADCAST}, and \cmd{SEND\_INCOMING}
are the three message types a client can receive \textbf{unsolicited}
(not as a direct reply to a command it just issued) --- a client
implementation must be able to tell these apart from ordinary replies
(see Section~8).

\section{Error Codes}

\begin{longtable}{@{}p{4.3cm}p{10.2cm}@{}}
\toprule
\textbf{Code} & \textbf{Sent when} \\
\midrule
\endhead
\cmd{BAD\_ARGS} & Wrong number of fields, or a non-integer where a size was expected. \\
\cmd{AUTH\_FAILED} & Username exists and the password doesn't match. \\
\cmd{USER\_ALREADY\_ONLINE} & That username already has an active session. \\
\cmd{NOT\_LOGGED\_IN} & Any command other than \cmd{LOGIN} sent before authenticating. \\
\cmd{ALREADY\_LOGGED\_IN} & \cmd{LOGIN} sent on a connection that already authenticated. \\
\cmd{UNKNOWN\_COMMAND} & Command name not recognized. \\
\cmd{USER\_NOT\_FOUND} & \cmd{MSG} or \cmd{SEND} target is not currently online. \\
\cmd{FILE\_NOT\_FOUND} & \cmd{DOWNLOAD} requested a file that doesn't exist in server storage. \\
\bottomrule
\end{longtable}

\section{Sequence Diagrams}

\subsection{Login}
\begin{verbatim}
Client                          Server
  |--- LOGIN|alice|secret ------->|
  |                                | check credentials, register session
  |<-------- LOGIN_OK|alice -------|
\end{verbatim}

\subsection{Private message}
\begin{verbatim}
Alice                Server                Bob
  |-- MSG|bob|hi --->|                      |
  |                   |-- MSG|alice|hi ---->|   (unsolicited push)
  |<---- MSG_OK ------|                      |
\end{verbatim}

\subsection{Upload / Download}
\begin{verbatim}
Client                          Server
  |-- UPLOAD|f.txt|1024 --------->|
  |<---------- READY -------------|
  |-- <1024 raw bytes> ---------->|
  |<------ UPLOAD_OK|f.txt -------|

Client                          Server
  |-- DOWNLOAD|f.txt ------------>|
  |<-- DOWNLOAD_OK|f.txt|1024 ----|
  |<----- <1024 raw bytes> -------|
\end{verbatim}

\subsection{Client-to-client SEND (relayed live through the server)}
\begin{verbatim}
Alice                    Server                    Bob
  |-- SEND|bob|f|1024 -->|                          |
  |<------ READY --------|                          |
  |-- <1024 raw bytes> ->|                          |
  |                       |-- SEND_INCOMING|f|alice|1024 -->|  (push)
  |                       |------- <1024 raw bytes> ------->|
  |<---- SEND_OK|f -------|                          |
\end{verbatim}

The server does not persist \cmd{SEND} payloads to disk --- it pulls the
bytes off the sender's socket and immediately relays them to the
target's socket. This means the target must be online at the time of
the \cmd{SEND} (offline delivery is out of scope for this assignment).

\section{``Active Users'' --- Design Note}

The brief's Server Responsibilities section asks the server to
``maintain a list of active users.'' That phrase is ambiguous between
(a)~users currently connected and (b)~every registered account. Both
readings are covered here rather than picking one:

\begin{itemize}[leftmargin=1.5em]
  \item \cmd{LIST} returns (a) --- who you can actually \cmd{MSG} right
        now, matching how the brief's own \cmd{LIST} example is used
        alongside \cmd{MSG}/\cmd{BROADCAST}.
  \item \cmd{ALLUSERS} returns (b) --- every account that has ever
        registered, each tagged \cmd{online}/\cmd{offline}, for a full
        directory view.
\end{itemize}

\section{Distinguishing Pushes from Replies (Client Design Note)}

Because a single TCP connection carries both the replies to commands the
client itself issued \textit{and} asynchronous pushes from other users,
the client's read loop must disambiguate by command name: \cmd{MSG},
\cmd{BROADCAST}, and \cmd{SEND\_INCOMING} are always pushes and are
handled immediately (printed, or --- for \cmd{SEND\_INCOMING} --- read
and saved to disk); every other command name is treated as the reply to
whatever request the client's main thread is currently waiting on.

\end{document}
